\section{原则四：优化 Token 效率}

优化上下文的\textbf{质量}很重要，但优化返回给 Agent 的上下文\textbf{数量}同样关键。

\subsection{控制响应大小的策略}

对于可能消耗大量上下文的 Tool 响应，建议实现以下机制的组合：

\begin{table}[H]
\centering
\caption{控制 Tool 响应大小的四种策略}
\begin{tabular}{p{0.18\textwidth}p{0.7\textwidth}}
\toprule
\textbf{策略} & \textbf{说明} \\
\midrule
分页 (Pagination) & 将大量结果分批返回，Agent 按需请求下一页 \\
\midrule
范围选择 (Range Selection) & 允许 Agent 指定返回结果的范围 \\
\midrule
过滤 (Filtering) & 允许 Agent 通过条件过滤只获取相关结果 \\
\midrule
截断 (Truncation) & 对超长响应自动截断，并附带有用的引导信息 \\
\bottomrule
\end{tabular}
\end{table}

\begin{importantbox}{Claude Code 的 Token 限制实践}
在 Claude Code 中，Anthropic 默认将 Tool 响应限制为 \textbf{25,000 tokens}。虽然他们预计 Agent 的有效上下文长度会随时间增长，但对 Token 高效工具的需求将始终存在。
\end{importantbox}

\subsection{引导 Agent 的高效行为}

如果选择截断响应，请确保附带有用的引导指令。可以：

\begin{itemize}[leftmargin=1.5em]
    \item 直接鼓励 Agent 采用更 Token 高效的策略，例如进行多次小型定向搜索，而非一次大范围搜索
    \item 在错误响应中清晰传达具体、可操作的改进建议，而非返回不透明的错误代码或调用栈
\end{itemize}

\begin{knowledgebox}{错误响应也是 Prompt Engineering}
Tool 的截断和错误响应本身就是一种 prompt engineering。它们可以引导 Agent 采用更 Token 高效的工具使用行为（如使用过滤器或分页），或展示正确格式化的工具输入示例。将错误响应视为与 Agent 沟通的机会，而非简单的失败通知。
\end{knowledgebox}

\subsection{本章小结}

Token 效率直接影响 Agent 的表现。通过分页、过滤、截断等策略控制响应大小，并利用截断和错误响应来引导 Agent 走向高效行为。Claude Code 的 25,000 token 默认限制是一个实用参考。


